% Report fragment. Requires amsmath, amsthm and theorem/lemma environments.
\section{A sparse scheduler that retains every finite proof word}
\label{sec:sparse-scheduler}

The source-level change is to visit every other fresh command code.
It preserves the numerical calculus, the AX, MP and Shift operations,
their stack representation, and their proofs of soundness. It changes the
infinite enumeration, rather than only its finite starting segment.
We first prove the arithmetic completeness of this change independently
of any proposed physical implementation.

Write
\[
 P(a,b)=\frac{(a+b)(a+b+1)}2+b.
\]
For a command $C$ and a continuation code $t$, the supplied command encoder is
\[
 \begin{array}{ll}
 e_{\mathsf S}(t)=P(0,t),&
 e_{\mathsf{MP}(u)}(t)=P(1,P(u,t)),\\[2pt]
 e_{\mathsf{AX}(c,x,y)}(t)=P(c+2,P(y,P(x,t))).
 \end{array}
\]
For a command word $W$, put $E([],t)=t$ and
$E(C::V,t)=e_C(E(V,t))+1$. If $W$ is nonempty, its fresh candidate code is
$F_W(t)=E(W,t)-1$. The literal decoder reads a fresh candidate without
subtracting one; it subtracts one from each positive continuation before
decoding that continuation. Consequently, for $t>0$, the word decoded from
$F_W(t)$ begins with exactly $W$, followed by the word decoded from the
residual pointer $t$. This is the supplied fresh/positive-pointer convention,
including the exceptional behavior of pointer zero.

Define the pairing depth
\[
 d(W)=\#\mathsf S(W)+2\#\mathsf{MP}(W)+3\#\mathsf{AX}(W).
\]
This counts pairings along the unique chain from the terminal pointer to
the fresh candidate; additions of one do not contribute to this count.

\begin{lemma}[Parity padding]\label{lem:parity-padding}
For every nonempty $W$ and every $t\geq0$, the difference
\[
             F_W(t+2^{d(W)})-F_W(t)
\]
is odd.
\end{lemma}
\begin{proof}
For $h=2^ku$, where $k\geq1$ and $u$ is odd, direct expansion gives
\[
 P(a,x+h)-P(a,x)
   =2^{k-1}u\bigl(2a+2x+h+3\bigr).
\]
The factor in parentheses is odd. Thus one pairing decreases the exact
power of two dividing the difference by one. Adding the same integer to
the two arguments leaves their difference unchanged. Starting from a
terminal-pointer difference $2^{d(W)}$ and proceeding through the
$d(W)$ pairings leaves an odd difference. All functions in the chain
are increasing, and $E(W,t)>0$ for nonempty $W$, so the final common
subtraction of one introduces no truncated-subtraction issue.
\end{proof}

\begin{theorem}[Completeness of a parity progression]
\label{thm:sparse-prefix}
For every $N_0\geq0$ and every nonempty finite command word $W$, there
is a $k\geq0$ such that the fresh word at candidate $N_0+2k$ has $W$
as a prefix. The candidate can be chosen primitive recursively in
$N_0$ and $W$, and PA proves the statement and the correctness of this
choice.
\end{theorem}
\begin{proof}
Take $t=N_0+1$. By Lemma~\ref{lem:parity-padding}, the two integers
$F_W(t)$ and $F_W(t+2^{d(W)})$ have opposite parity. Select the one
congruent to $N_0$ modulo two. Pairing satisfies $P(a,b)\geq b$,
so both candidates are at least $t>N_0$. The selected candidate is
therefore $N_0+2k$ for a natural $k$. Both choices have the required
word prefix by the encoding identity. Code formation, exponentiation,
parity testing, and the selected division by two are primitive recursive.
The proof uses arithmetic identities and induction over a finite word,
all available in PA.
\end{proof}

\subsection{The exact register-program edit}

The nine registers are ordered $(N,T,W,A,D,H,X,Y,L)$. In the primary
PA194 source, the relevant shared trees are
\[
\begin{aligned}
 \mathsf{node27}&=(\mathsf{node25},\mathsf{node26}),\\
 \mathsf{node25}&=(\operatorname{dec}N,
        (\operatorname{inc}L,(\operatorname{inc}X,\operatorname{continue}3))),\\
 \mathsf{node26}&=(\operatorname{inc}N,
        (\operatorname{dec}X,\operatorname{continue}2)).
\end{aligned}
\]
The first loop drains $N$ into $L$ and $X$. The second loop restores
$X$ into $N$ and performs one additional increment. To advance the fresh
counter by two, replace only the first line by
\[
 \mathsf{node27}'=(\mathsf{node25},
          (\operatorname{inc}N,\mathsf{node26})).
\]
The additional increment is outside the restoring loop. Duplicating
its increment inside that loop would instead change $N$ to $2(N+1)$
and would require a different completeness argument.

At the resolved instruction level, insert an increment of register zero
at ordinal 31. Identify each old ordinal $i$ with $i$ for $i<31$ and
with $i+1$ otherwise. All old instructions and successors are preserved
under this identification except the zero branch of the decrement at
ordinal 27: that branch enters the new ordinal 31, which then enters
ordinal 32, the relocated old increment. The resulting source has
3,674 primitive instructions and 4,781 branch-sensitive contracts.
Its maximum source-path length is 36.

Induction on the drained register and then on the restored register gives
an exact fresh-fetch contract
\[
 (N,L,X)=(n,0,0)\quad\longmapsto\quad(n+2,n,0)
\]
in $7n+5$ source steps, from ordinal 25 to the relocated unpairing
entry at ordinal 35. All other registers are unchanged. For $L=l+1>0$,
the unchanged positive-pointer branch takes two source steps and yields
$(N,L,X)=(n,l,0)$. In particular, the added increment is executed exactly
once per fresh candidate, and never once per continued command.

\paragraph{Joining an intermediate physical-startup endpoint.}
The verified source trace supplied with this study starts at ordinal 40
with logical vector $(4,0,0,0,0,0,0,0,1)$. It reaches ordinal 25 with
that same vector after 29 source instructions. The positive continuation
$L=1$ then decodes command code zero, which is Shift. From the zero proof
store this Shift is harmless; after 19 further source instructions the
machine returns to ordinal 25 with vector $(4,0,0,0,0,0,0,0,0)$.
The delivered trace lists all 49 configurations of this 48-instruction
prefix and contains no halt. Thus a literal blank-start certificate ending
at the former frame joins to the required clean seed with $N_0=4$.
The change of physical register locations, when used by the physical
construction, has no effect on this source calculation.

\subsection{PA consistency calibration of the changed source}

Let $R_2$ denote this modified source, with a clean proof store, scratch
registers zero and arbitrary initial fresh counter $N_0$. As in the
supplied calibration, every visible theorem code and every canonical
stack entry is annotated by a PA proof of its arithmetical translation.
The scheduler change acts only on the fresh-code counter. It does not
change a command, its logical output, its halt test, or any proof-store
update. The same syntax-level annotation therefore transforms every
finite halting run of $R_2$ into a PA refutation.

For the other direction, apply the supplied primitive-recursive
refutation compiler to a PA refutation. It supplies a nonempty command
word $W$ that halts from \emph{every} incoming command store. This
arbitrary-store property matters: the scheduler retains its proof store
between candidates. The initial AX command of the compiler clears the
caches; the subsequent proof-producing word works over an arbitrary
saved stack. For the separate-head calculus the ten-entry stack offsets
must remain those of that calculus; the scheduler change does not alter
them.

Theorem~\ref{thm:sparse-prefix} provides a candidate $N_0+2k$ whose decoded
word begins with $W$. Each preceding candidate either halts, already
satisfying the conclusion, or completes after finitely many commands.
A positive continuation strictly decreases, since it is decoded from
$L-1$ and the next component is at most $L-1$. Therefore every finite
candidate completes unless it halts. Induction through the first $k$
candidates reaches the selected one unless an earlier halt occurs.
The selected candidate halts during its prefix $W$, regardless of the
store left by the preceding candidates. Halting is absorbing, so the
arbitrary residual suffix is immaterial.

All these operations are primitive-recursive transformations of finite
codes. Their loop bounds and trace concatenations are proved by ordinary
arithmetic induction; there is no appeal to standard-model soundness or
to a reflection principle for PA. Thus, writing $H_{R_2}$ for finite
halting from the specified clean seed,
\[
       \mathrm{PA}\vdash
       H_{R_2}\;\longleftrightarrow\;
       \exists p\,\operatorname{Prf}_{\mathrm{PA}}
                    (p,\ulcorner0=1\urcorner).
\]
Any literal binary machine whose exact blank startup reaches this seed,
and whose positive finite simulations implement all modified source
instructions, inherits the desired PA-provable nonhalting biconditional.
The source theorem itself does not assert a physical state count.
